\section*{Leitidee}

Vektorr�ume beschreiben mathematische Objekte, die sich addieren und mit
Skalaren multiplizieren lassen. Der n�chste Schritt besteht darin,
Abbildungen zwischen solchen R�umen zu untersuchen, welche diese lineare
Struktur erhalten.

Solche \textbf{linearen Abbildungen} lassen sich in endlichdimensionalen
Vektorr�umen mithilfe von \textbf{Matrizen} darstellen. Eine Matrix ist damit
nicht nur eine rechteckige Anordnung von Zahlen, sondern beschreibt, wie eine
lineare Abbildung auf Vektoren wirkt. Matrixoperationen erhalten dadurch eine
inhaltliche Bedeutung: Sie erm�glichen es, lineare Transformationen
algebraisch zu beschreiben und miteinander zu verkn�pfen.

Gleichzeitig bilden Matrizen die nat�rliche Sprache f�r
\textbf{lineare Gleichungssysteme}. Mehrere lineare Gleichungen k�nnen in
einer einzigen Matrixgleichung zusammengefasst werden. Dadurch entstehen
systematische Verfahren, mit denen sich L�sbarkeit, L�sungsr�ume und
Abh�ngigkeiten untersuchen lassen.

Mit der \textbf{Determinante} und der \textbf{inversen Matrix} kommen
schlie�lich Gr��en und Operationen hinzu, die wesentliche Eigenschaften
linearer Abbildungen sichtbar machen. Insbesondere l�sst sich untersuchen,
ob eine lineare Transformation umkehrbar ist und ob ein lineares
Gleichungssystem eine eindeutig bestimmte L�sung besitzt.

Die dabei entwickelten Begriffe reichen weit �ber das Rechnen mit
Zahlenspalten und Matrizen hinaus. In sp�teren Kapiteln wird die lineare
Struktur auf Funktionen �bertragen. Linearkombinationen, lineare
Abbildungen und Zerlegungen bilden dort unter anderem die Grundlage f�r
Approximationen und die Fourier-Analyse.

